\documentclass[12pt]{article}
\usepackage{amsmath}
\usepackage{amsfonts}
\usepackage{float}
\usepackage{fancyhdr}
\usepackage{graphicx}
\usepackage{setspace}
\usepackage[colorlinks=true,linkcolor=blue, citecolor=red]{hyperref}
\usepackage{url}
\usepackage[top=.75in, left=.75in, right=.75in, bottom=1in]{geometry}

% For algorithm
\usepackage{algorithm}
\usepackage{algpseudocode}

% ============ CODE ============
\usepackage{listings}
\usepackage{xcolor}
\definecolor{codegreen}{rgb}{0,0.6,0}
\definecolor{codegray}{rgb}{0.5,0.5,0.5}
\definecolor{codepurple}{rgb}{0.58,0,0.82}
\definecolor{backcolour}{rgb}{0.95,0.95,0.92}

% Styling for the code.
\lstdefinestyle{mystyle}{
    backgroundcolor=\color{backcolour},   
    commentstyle=\color{codegreen},
    keywordstyle=\color{magenta},
    numberstyle=\tiny\color{codegray},
    stringstyle=\color{codepurple},
    basicstyle=\ttfamily\footnotesize,
    breakatwhitespace=false,         
    breaklines=true,                 
    captionpos=b,                    
    keepspaces=true,                 
    numbers=left,                    
    numbersep=5pt,                  
    showspaces=false,                
    showstringspaces=false,
    showtabs=false,                  
    tabsize=2
}
\lstset{style=mystyle}

% Table packages and Unicode / Font support
\usepackage{array}
\usepackage{longtable}
\usepackage{enumitem}
\usepackage{booktabs}
\usepackage[utf8]{inputenc}
\DeclareUnicodeCharacter{2264}{\ensuremath{\le}}
\DeclareUnicodeCharacter{2194}{\ensuremath{\leftrightarrow}}
\usepackage[T1]{fontenc}
\usepackage{lmodern}
\graphicspath{{../../../../templates/}{../../prototype/}{figures/}{./}}
\setlength{\emergencystretch}{3em}
\setlist[enumerate]{label=\arabic*.,leftmargin=*}
\setlist[itemize]{leftmargin=*}

% ============ HEADER AND FOOTER ============
\setlength{\headheight}{29.43912pt}

\pagestyle{fancy}
\fancyfoot{}
\fancyfoot[R]{Page \thepage}

\setlength{\parindent}{0pt}
\renewcommand{\baselinestretch}{1.5}

\newcommand{\coursename}{Introduction of Software Engineering}
\newcommand{\reportname}{Studitopia}
\newcommand{\reporttitle}{\href{https://github.com/RoundFloor/Studitopia.git}{Studitopia}}

\newcommand{\studentname}{%
Nguyen Thanh An (24125024)\\
Nguyen Chi Bao (24125025)\\
Tran Nguyen Khanh Binh (24125026)\\
Tran Canh Anh Tuan (24125048)\\
Nguyen Le Quoc Huy (24125100)
}
\newcommand{\teachername}{Ho Tuan Thanh, Nguyen Thi Minh Tuyen}

% Header
\lhead{\reporttitle}
\rhead{
University of Science - VNUHCM\\
\coursename
}

% Footer
\newcommand{\leftfooter}{}
\lfoot{\leftfooter}

% ============ DOCUMENT ============
\begin{document}

\pagenumbering{roman}
\begin{titlepage}
\newcommand{\HRule}{\rule{\linewidth}{0.5mm}}
\centering

\textsc{\LARGE VNU-HCM}\\[0.5cm]
\textsc{\Large University of Science}\\[0.5cm]
\textsc{\large Faculty of Information Technology}\\[0.5cm]
\textsc{Computer Science}\\[0.5cm]

\HRule \\[0.4cm]
{ 
\huge{\bfseries{\reporttitle}}\\[0.5cm]
\large{\bfseries{UC-24 --- View My Storage}}
}\\[0.4cm]
\HRule \\[0.5cm]

\textbf{\large Subject: \coursename}\\[0.5cm]

\begin{minipage}[t]{0.55\textwidth}
\begin{flushleft} \large
\emph{Students:}\\
\studentname
\end{flushleft}
\end{minipage}
~
\begin{minipage}[t]{0.4\textwidth}
\begin{flushright} \large
\emph{Supervising Teachers:} \\
\teachername
\end{flushright}
\end{minipage}\\[1cm]

{\large \today}\\[1cm]

\includegraphics[scale=.20]{img/hcmus-logo.png}\\[1cm] 

\vfill
\end{titlepage}	

\tableofcontents
\pagebreak

\pagenumbering{arabic}
\setcounter{page}{1}

\hypertarget{sec:uc-24-view-my-storage}{}
\addcontentsline{toc}{section}{UC-24 --- View My Storage}
\section*{UC-24 --- View My Storage}
\par\medskip

\textbf{Version:} \texttt{1.0}
\par\medskip

\textbf{Performed by:} [Nguyen Le Quoc Huy] | \textbf{Edited by:} [Nguyen Le Quoc Huy]
\par\medskip

\begingroup
\small
\renewcommand{\baselinestretch}{1.0}\selectfont
\renewcommand{\arraystretch}{1.15}
\setlength{\tabcolsep}{6pt}
\begin{longtable}{|>{\raggedright\arraybackslash}p{0.09\textwidth}|>{\raggedright\arraybackslash}p{0.84\textwidth}|}
\hline
\textbf{Field} & \textbf{Value} \\ \hline
\endhead
Use Case ID & UC-24 \\ \hline
Name & View My Storage \\ \hline
Primary Actor & University Student \\ \hline
Supporting Actors & None \\ \hline
Feature Group & Learning Resource Library \\ \hline
Priority & Not explicitly stated \\ \hline
Related User Stories & US-28 \\ \hline
Related Requirements & FR-28.1 to FR-28.3; AC-28.1 to AC-28.3; FR-27.1; FR-27.5; NFR-PERF-1; NFR-SCAL-4; NFR-SEC-1; NFR-SEC-6; NFR-PRIV-5; NFR-USE-4 \\ \hline
Source Documents & \texttt{docs/requirements/proposals/userstory.md}; \texttt{docs/requirements/proposals/functional\_requirement.md}; \texttt{docs/requirements/proposals/library.md}; \texttt{docs/requirements/proposals/create\_post.md}; \texttt{docs/requirements/proposals/API.md}; \texttt{docs/requirements/proposals/response\_code.md}; \texttt{docs/requirements/proposals/non\_functional\_requirement.md}; \texttt{docs/requirements/proposals/feature\_flow.md}; \texttt{docs/requirements/proposals/screen\_flow.md} \\ \hline
\end{longtable}
\endgroup
\par\medskip

\hypertarget{sec:part-a-use-case-model}{}
\addcontentsline{toc}{subsection}{Part A --- Use Case Model}
\subsection*{Part A --- Use Case Model}
\par\medskip

\hypertarget{sec:a1-focused-use-case-diagram}{}
\addcontentsline{toc}{subsubsection}{A1. Focused Use Case Diagram}
\subsubsection*{A1. Focused Use Case Diagram}
\par\medskip

\begin{center}
\includegraphics[width=0.85\textwidth,height=0.75\textheight,keepaspectratio]{"figures/UC-24-view-my-storage_1.png"}
\end{center}
\par\medskip

\textbf{Diagram notation.} Person shape = human actor; cloud = supporting system; blue = selected use case; solid line = association; labelled solid arrow = generalization; dashed arrow = include or extend.
\par\medskip

\hypertarget{sec:a2-actor-register}{}
\addcontentsline{toc}{subsubsection}{A2. Actor Register}
\subsubsection*{A2. Actor Register}
\par\medskip

\begingroup
\small
\renewcommand{\baselinestretch}{1.0}\selectfont
\renewcommand{\arraystretch}{1.15}
\setlength{\tabcolsep}{6pt}
\begin{longtable}{|>{\raggedright\arraybackslash}p{0.13\textwidth}|>{\raggedright\arraybackslash}p{0.09\textwidth}|>{\raggedright\arraybackslash}p{0.69\textwidth}|}
\hline
\textbf{Actor} & \textbf{Type} & \textbf{Description} \\ \hline
\endhead
University Student & Primary & Authenticated resource owner who reviews published resources and their view and access counts. \\ \hline
\end{longtable}
\endgroup
\par\medskip

\hypertarget{sec:a3-relationships}{}
\addcontentsline{toc}{subsubsection}{A3. Relationships}
\subsubsection*{A3. Relationships}
\par\medskip

\hypertarget{sec:a3-1-associations}{}
\paragraph*{A3.1 Associations}
\par\medskip

\begingroup
\small
\renewcommand{\baselinestretch}{1.0}\selectfont
\renewcommand{\arraystretch}{1.15}
\setlength{\tabcolsep}{6pt}
\begin{longtable}{|>{\raggedright\arraybackslash}p{0.16\textwidth}|>{\raggedright\arraybackslash}p{0.18\textwidth}|>{\raggedright\arraybackslash}p{0.10\textwidth}|>{\raggedright\arraybackslash}p{0.47\textwidth}|}
\hline
\textbf{Actor} & \textbf{Use Case} & \textbf{Type} & \textbf{Description} \\ \hline
\endhead
University Student & UC-24 View My Storage & Primary & Opens a personal storage view of published resources. \\ \hline
\end{longtable}
\endgroup
\par\medskip

\hypertarget{sec:a3-2-generalization-and-specialization}{}
\paragraph*{A3.2 Generalization and Specialization}
\par\medskip

\begingroup
\small
\renewcommand{\baselinestretch}{1.0}\selectfont
\renewcommand{\arraystretch}{1.15}
\setlength{\tabcolsep}{6pt}
\begin{longtable}{|>{\raggedright\arraybackslash}p{0.09\textwidth}|>{\raggedright\arraybackslash}p{0.09\textwidth}|>{\raggedright\arraybackslash}p{0.26\textwidth}|>{\raggedright\arraybackslash}p{0.47\textwidth}|}
\hline
\textbf{Child} & \textbf{Parent} & \textbf{Type} & \textbf{Inherited Behavior or Association} \\ \hline
\endhead
None & None & No generalization or specialization & The proposals name only the University Student actor for US-28. \\ \hline
\end{longtable}
\endgroup
\par\medskip

\hypertarget{sec:a3-3-include}{}
\paragraph*{A3.3 Include}
\par\medskip

\begingroup
\small
\renewcommand{\baselinestretch}{1.0}\selectfont
\renewcommand{\arraystretch}{1.15}
\setlength{\tabcolsep}{6pt}
\begin{longtable}{|>{\raggedright\arraybackslash}p{0.23\textwidth}|>{\raggedright\arraybackslash}p{0.30\textwidth}|>{\raggedright\arraybackslash}p{0.40\textwidth}|}
\hline
\textbf{Base Use Case} & \textbf{Included Use Case} & \textbf{Meaning} \\ \hline
\endhead
None & None & No include relationship & The proposals do not name an included use case for US-28. \\ \hline
\end{longtable}
\endgroup
\par\medskip

\hypertarget{sec:a3-4-extend}{}
\paragraph*{A3.4 Extend}
\par\medskip

\begingroup
\small
\renewcommand{\baselinestretch}{1.0}\selectfont
\renewcommand{\arraystretch}{1.15}
\setlength{\tabcolsep}{6pt}
\begin{longtable}{|>{\raggedright\arraybackslash}p{0.15\textwidth}|>{\raggedright\arraybackslash}p{0.11\textwidth}|>{\raggedright\arraybackslash}p{0.18\textwidth}|>{\raggedright\arraybackslash}p{0.47\textwidth}|}
\hline
\textbf{Extending Use Case} & \textbf{Base Use Case} & \textbf{Condition} & \textbf{Added Behavior} \\ \hline
\endhead
None & None & No extend relationship & The proposals do not name an extending use case for US-28. \\ \hline
\end{longtable}
\endgroup
\par\medskip

\hypertarget{sec:a4-model-notes}{}
\addcontentsline{toc}{subsubsection}{A4. Model Notes}
\subsubsection*{A4. Model Notes}
\par\medskip

\begin{enumerate}
\item The scope is the owner's personal storage view, status and publication-date display, and view/access statistics.
\item Publishing resources is outside this use case; deletion is handled in US-29 and existing UC-06 covers the publication flow.
\item The status values shown in the storage table are exactly \texttt{Published}, \texttt{Restricted}, and \texttt{Unavailable}.
\item No distinct included, extending, or specialized use case is named for US-28, so no I-xx, E-xx, or G-xx relationship ID is introduced.
\item The exact API endpoint, route, pagination limit, and priority are not specified in proposals.
\end{enumerate}
\par\medskip

\hypertarget{sec:part-b-use-case-specification}{}
\addcontentsline{toc}{subsection}{Part B --- Use Case Specification}
\subsection*{Part B --- Use Case Specification}
\par\medskip

\hypertarget{sec:b1-brief-description}{}
\addcontentsline{toc}{subsubsection}{B1. Brief Description}
\subsubsection*{B1. Brief Description}
\par\medskip

A university student views personal storage so that the student can track how many views and accesses published resources receive. The system provides a personal storage view for published resources, displays view and access counts for each resource, and displays publication status and date.
\par\medskip

\texttt{library.md} recommends presenting this as \textbf{My Storage}, with summary totals and a table of owned resources. These UI details are recommendations; the confirmed requirements are \texttt{AC-28.1} to \texttt{AC-28.3} and \texttt{FR-28.1} to \texttt{FR-28.3}.
\par\medskip

\hypertarget{sec:b2-preconditions}{}
\addcontentsline{toc}{subsubsection}{B2. Preconditions}
\subsubsection*{B2. Preconditions}
\par\medskip

\begin{enumerate}
\item The student is authenticated before accessing protected user resources (\texttt{NFR-SEC-1}).
\item The Learning Resource Library is reachable through \texttt{Home -$>$ Library Group} in \texttt{screen\_flow.md}.
\item The student has published resources through the library or material-post path (\texttt{FR-27.1}, \texttt{FR-15.6}, \texttt{FR-27.4}).
\item View, access, status, and publication-date data are available for the owned resources (\texttt{FR-28.2}, \texttt{FR-28.3}).
\item Authorization checks restrict storage data to the resource owner (\texttt{NFR-SEC-6}).
\end{enumerate}
\par\medskip

\hypertarget{sec:b3-postconditions}{}
\addcontentsline{toc}{subsubsection}{B3. Postconditions}
\subsubsection*{B3. Postconditions}
\par\medskip

\begin{enumerate}
\item \textbf{Success:} The student's published resources are listed in a personal storage view.
\item \textbf{Success:} Each listed resource shows view count, access count, publication status, and publication date.
\item \textbf{Success:} The student can open a resource detail view or proceed to the delete action where available.
\item \textbf{Failure:} No other student's storage data is exposed, and the specified empty-state or load-error feedback is shown.
\end{enumerate}
\par\medskip

\hypertarget{sec:b4-main-success-scenario-basic-flow}{}
\addcontentsline{toc}{subsubsection}{B4. Main Success Scenario (Basic Flow)}
\subsubsection*{B4. Main Success Scenario (Basic Flow)}
\par\medskip

\begin{enumerate}
\item The student opens the library tab bar and selects \textbf{My Storage} through \texttt{MyStorageTab}.
\item The system displays \texttt{MyStoragePanel} with heading \textbf{My Storage}.
\item The system displays \texttt{MyStorageSummaryStrip} with total published resources, total views, and total accesses/downloads.
\item The system displays \texttt{MyPublishedResourcesTable} with one row per published resource.
\item Each row shows the resource title and type through \texttt{MyResourceTitleCell}.
\item Each row shows a \texttt{status} badge with exactly \texttt{Published}, \texttt{Restricted}, or \texttt{Unavailable}.
\item Each row shows the publication date and time through \texttt{MyResourcePublishedDateCell} (\texttt{FR-28.3}, \texttt{AC-28.3}).
\item Each row shows the number of views through \texttt{MyResourceViewsCell} and the number of downloads/accesses through \texttt{MyResourceAccessesCell} (\texttt{FR-28.2}, \texttt{AC-28.2}).
\item Each row shows the privacy setting through \texttt{MyResourcePrivacyCell}.
\item The student selects \textbf{View} through \texttt{MyResourceOpenButton} to open the resource detail view, or remains in My Storage.
\end{enumerate}
\par\medskip

\hypertarget{sec:b5-alternative-flows}{}
\addcontentsline{toc}{subsubsection}{B5. Alternative Flows}
\subsubsection*{B5. Alternative Flows}
\par\medskip

\hypertarget{sec:a1-no-published-resources}{}
\paragraph*{A1. No published resources}
\par\medskip

The student has no published resources. The system displays the exact empty-state message “You have not published any resources yet.”
\par\medskip

\hypertarget{sec:a2-restricted-or-unavailable-resource}{}
\paragraph*{A2. Restricted or unavailable resource}
\par\medskip

A listed resource has status \texttt{Restricted} or \texttt{Unavailable}. The system still shows the owner's status badge and statistics, while resource access follows the owner's privacy setting (\texttt{FR-27.5}, \texttt{NFR-PRIV-5}).
\par\medskip

\hypertarget{sec:a3-open-resource-detail}{}
\paragraph*{A3. Open resource detail}
\par\medskip

The student selects \textbf{View}. The system opens the resource detail page with full metadata, preview where supported, rating, access information, views, and accesses.
\par\medskip

\hypertarget{sec:a4-retry-after-an-error}{}
\paragraph*{A4. Retry after an error}
\par\medskip

The student selects \textbf{Retry} through \texttt{StorageRetryButton}. The system reloads the storage data.
\par\medskip

\hypertarget{sec:b6-exception-flows}{}
\addcontentsline{toc}{subsubsection}{B6. Exception Flows}
\subsubsection*{B6. Exception Flows}
\par\medskip

\hypertarget{sec:e1-storage-data-load-failure}{}
\paragraph*{E1. Storage data load failure}
\par\medskip

The storage view cannot be loaded. The system shows “Unable to load resources.” with \textbf{Retry} through \texttt{LibraryErrorBanner}.
\par\medskip

\hypertarget{sec:e2-statistics-unavailable}{}
\paragraph*{E2. Statistics unavailable}
\par\medskip

View or access counts are unavailable for a resource. The system does not present an invented number; the exact missing-statistics message is Not specified in proposals.
\par\medskip

\hypertarget{sec:e3-unauthorized-storage-access}{}
\paragraph*{E3. Unauthorized storage access}
\par\medskip

A user who is not the owner attempts to open another student's storage view. The system denies access through server-side authorization checks (\texttt{NFR-SEC-6}); the exact storage-denial message is Not specified in proposals.
\par\medskip

\hypertarget{sec:b7-business-rules}{}
\addcontentsline{toc}{subsubsection}{B7. Business Rules}
\subsubsection*{B7. Business Rules}
\par\medskip

\begin{enumerate}
\item The system provides a personal storage view for published resources (\texttt{FR-28.1}, \texttt{AC-28.1}).
\item The system displays the number of views and accesses for each published resource (\texttt{FR-28.2}, \texttt{AC-28.2}).
\item The system displays the publication status and date of each resource (\texttt{FR-28.3}, \texttt{AC-28.3}).
\item Storage status values are exactly \texttt{Published}, \texttt{Restricted}, and \texttt{Unavailable}.
\item Privacy settings control access to documents and learning resources the student shared (\texttt{FR-27.5}, \texttt{NFR-PRIV-5}).
\item Only the owner may view personal storage data (\texttt{NFR-SEC-1}, \texttt{NFR-SEC-6}).
\end{enumerate}
\par\medskip

\hypertarget{sec:b8-special-requirements}{}
\addcontentsline{toc}{subsubsection}{B8. Special Requirements}
\subsubsection*{B8. Special Requirements}
\par\medskip

\begingroup
\small
\renewcommand{\baselinestretch}{1.0}\selectfont
\renewcommand{\arraystretch}{1.15}
\setlength{\tabcolsep}{6pt}
\begin{longtable}{|>{\raggedright\arraybackslash}p{0.10\textwidth}|>{\raggedright\arraybackslash}p{0.83\textwidth}|}
\hline
\textbf{ID} & \textbf{Requirement} \\ \hline
\endhead
NFR-PERF-1 & The system shall provide responsive user interactions with a response time of ≤ 2 seconds for standard user actions. \\ \hline
NFR-SCAL-4 & The system architecture shall support at least 100,000 uploaded documents and learning resources. \\ \hline
NFR-SEC-1 & The system shall require authentication before accessing protected user resources. \\ \hline
NFR-SEC-6 & The system shall prevent unauthorized access through server-side authorization checks. \\ \hline
NFR-PRIV-5 & Users shall be able to control access to documents and learning resources they share. \\ \hline
NFR-USE-4 & Primary interface actions shall use clear labels and recognizable controls. \\ \hline
Design constraint & Status values, table cells, summary cards, empty state, and labels follow \texttt{library.md} as UI recommendations; FR/AC/NFR entries are confirmed requirements. \\ \hline
\end{longtable}
\endgroup
\par\medskip

\hypertarget{sec:b9-extension-points}{}
\addcontentsline{toc}{subsubsection}{B9. Extension Points}
\subsubsection*{B9. Extension Points}
\par\medskip

\begingroup
\small
\renewcommand{\baselinestretch}{1.0}\selectfont
\renewcommand{\arraystretch}{1.15}
\setlength{\tabcolsep}{6pt}
\begin{longtable}{|>{\raggedright\arraybackslash}p{0.41\textwidth}|>{\raggedright\arraybackslash}p{0.52\textwidth}|}
\hline
\textbf{Extension Point} & \textbf{Location in Flow of Events} \\ \hline
\endhead
Empty storage & Before basic flow step 4; alternative A1 \\ \hline
Restricted or unavailable status & Basic flow step 6; alternative A2 \\ \hline
Resource detail access & Basic flow step 10; alternative A3 \\ \hline
Storage retry & After exception E1; alternative A4 \\ \hline
Storage load failure & Basic flow steps 2-4; exception E1 \\ \hline
Unauthorized access & Before basic flow step 2; exception E3 \\ \hline
\end{longtable}
\endgroup
\par\medskip

\hypertarget{sec:b10-interface-and-data-notes}{}
\addcontentsline{toc}{subsubsection}{B10. Interface and Data Notes}
\subsubsection*{B10. Interface and Data Notes}
\par\medskip

\begin{itemize}
\item \textbf{Screen and flow path (confirmed):} \texttt{screen\_flow.md} defines \texttt{Home -$>$ Library Group}. \texttt{feature\_flow.md} defines \texttt{Library -$>$ Material -$>$ Upload / Download}; \texttt{userstory.md} section 8.4 defines the personal storage requirement.
\item \textbf{Route (not specified):} Route is Not specified in proposals; \texttt{library.md} presents \texttt{/library/my-storage} as a suggested state only.
\item \textbf{Recommended UI:} \texttt{MyStorageTab} label \textbf{My Storage}, \texttt{MyStoragePanel} heading \textbf{My Storage}, \texttt{MyStorageSummaryStrip}, \texttt{MyPublishedResourcesTable}, \texttt{MyResourceStatusCell}, \texttt{MyResourcePublishedDateCell}, \texttt{MyResourceViewsCell}, \texttt{MyResourceAccessesCell}, \texttt{MyResourcePrivacyCell}, \texttt{MyResourceOpenButton}, \texttt{DeleteMyResourceButton}, and \texttt{StorageRetryButton}.
\item \textbf{Confirmed statistics and status:} \texttt{FR-28.2} / \texttt{AC-28.2} require views and accesses; \texttt{FR-28.3} / \texttt{AC-28.3} require publication status and date. The exact status values are \texttt{Published}, \texttt{Restricted}, and \texttt{Unavailable}.
\item \textbf{Exact messages and labels:} \textbf{My Storage}; \textbf{View}; \textbf{Delete}; \textbf{Retry}; “You have not published any resources yet.”; “Unable to load resources.” with \textbf{Retry}. No exact success message is specified for opening storage.
\item \textbf{Data fields (confirmed basis plus recommended names):} \texttt{resourceId}, \texttt{title}, \texttt{resourceType}, \texttt{privacy}, \texttt{views}, \texttt{accesses}, \texttt{status}, \texttt{publishedAt}, \texttt{rating}, \texttt{reviewCount}, and \texttt{sourceForumPostId} are named in \texttt{library.md}; only owner-visible data is shown in the storage table.
\item \textbf{API and codes:} Library storage endpoint, statuses, and error codes are Not specified in proposals. \texttt{API.md} has no standalone library storage endpoint, and \texttt{response\_code.md} defines no library-specific error code.
\item \textbf{Pagination and limits:} Row limits, pagination, and refresh interval are Not specified in proposals.
\item \textbf{Notification and privacy rules:} No story-specific notification is stated. Storage access must respect owner privacy settings (\texttt{FR-27.5}, \texttt{NFR-PRIV-5}) and server-side authorization (\texttt{NFR-SEC-6}).
\end{itemize}
\par\medskip

\hypertarget{sec:b11-source-alignment-notes}{}
\addcontentsline{toc}{subsubsection}{B11. Source Alignment Notes}
\subsubsection*{B11. Source Alignment Notes}
\par\medskip

\begin{itemize}
\item This use case is derived from US-28 in \texttt{userstory.md} section 8.4 and from \texttt{FR-28.1} to \texttt{FR-28.3} and \texttt{AC-28.1} to \texttt{AC-28.3} in the proposals.
\item \texttt{library.md} supplies recommended component names, status values, and exact UI messages; these are UI recommendations, not new FRs.
\item \texttt{API.md} has no standalone library storage endpoint, so no storage route, status code, or error code is invented.
\item Overlap with existing \texttt{UC-06-share-learning-resource.md}: UC-06 publishes a resource and directs the owner to the storage view; UC-24 defines what that storage view displays. Deletion is covered by US-29 rather than this use case.
\item No I-xx, E-xx, or G-xx model relationship ID is introduced; flow headings such as A2 and E2 are specification flow labels only.
\end{itemize}
\par\medskip

\end{document}
